\begin{proof}
Let  

\[
\mathbf a=(a_{1},\dots ,a_{n}),\qquad 
\mathbf b=(b_{1},\dots ,b_{n})
\]

be vectors in \(\mathbb{R}^{n}\).  
Their inner product and squared Euclidean norms are  

\[
\langle\mathbf a,\mathbf b\rangle=\sum_{i=1}^{n}a_{i}b_{i},\qquad 
\|\mathbf a\|^{2}= \sum_{i=1}^{n}a_{i}^{2},\qquad 
\|\mathbf b\|^{2}= \sum_{i=1}^{n}b_{i}^{2}.
\]

\medskip
\textbf{1. A non‑negative quadratic.}  
Define  

\[
f(t)=\sum_{i=1}^{n}(a_{i}t-b_{i})^{2},\qquad t\in\mathbb R .
\]

Each summand is a square, hence \(f(t)\ge 0\) for every real \(t\).

\medskip
\textbf{2. Expansion of \(f(t)\).}  
Using \((x-y)^{2}=x^{2}-2xy+y^{2}\) and linearity of the sum,

\begin{align*}
f(t)
&=\sum_{i=1}^{n}\bigl(a_{i}^{2}t^{2}-2a_{i}b_{i}t+b_{i}^{2}\bigr)\\
&=\Bigl(\sum_{i=1}^{n}a_{i}^{2}\Bigr)t^{2}
   -2\Bigl(\sum_{i=1}^{n}a_{i}b_{i}\Bigr)t
   +\sum_{i=1}^{n}b_{i}^{2}.
\end{align*}

Thus  

\[
f(t)=At^{2}+Bt+C,\qquad 
A=\sum_{i=1}^{n}a_{i}^{2},\;
B=-2\sum_{i=1}^{n}a_{i}b_{i},\;
C=\sum_{i=1}^{n}b_{i}^{2}.
\]

\medskip
\textbf{3. Discriminant condition.}  
If a quadratic \(At^{2}+Bt+C\) is non‑negative for all real \(t\) and \(A>0\), its discriminant satisfies  

\[
\Delta = B^{2}-4AC\le 0.
\]

(The case \(A=0\) occurs only when \(\mathbf a=\mathbf 0\); then \(f(t)=\sum b_{i}^{2}\ge0\) and the desired inequality is trivial.)

Substituting the expressions for \(A,B,C\),

\begin{align*}
\Delta
&= \bigl(-2\sum_{i=1}^{n}a_{i}b_{i}\bigr)^{2}
   -4\Bigl(\sum_{i=1}^{n}a_{i}^{2}\Bigr)\Bigl(\sum_{i=1}^{n}b_{i}^{2}\Bigr)\\
&= 4\Bigl(\sum_{i=1}^{n}a_{i}b_{i}\Bigr)^{2}
   -4\Bigl(\sum_{i=1}^{n}a_{i}^{2}\Bigr)\Bigl(\sum_{i=1}^{n}b_{i}^{2}\Bigr)\le 0 .
\end{align*}

Dividing by \(4\) gives the **Cauchy–Schwarz inequality**

\[
\boxed{\left(\sum_{i=1}^{n}a_{i}b_{i}\right)^{2}
      \le \left(\sum_{i=1}^{n}a_{i}^{2}\right)
         \left(\sum_{i=1}^{n}b_{i}^{2}\right)} .
\]

\medskip
\textbf{4. Equality condition.}  
Equality holds exactly when \(\Delta=0\), i.e. when the quadratic \(f(t)\) has a single real root.  
If \(\Delta=0\), there exists \(\lambda\in\mathbb R\) such that the root satisfies  

\[
a_{i}\lambda-b_{i}=0\qquad\text{for every }i,
\]

or equivalently  

\[
b_{i}= \lambda\,a_{i}\quad\text{for all }i .
\]

Hence the vectors are linearly dependent (one is a scalar multiple of the other).  
Conversely, if \(\mathbf b=\lambda\mathbf a\) for some \(\lambda\in\mathbb R\), then  

\[
f(t)=\sum_{i=1}^{n}\bigl(a_{i}t-\lambda a_{i}\bigr)^{2}
    =(t-\lambda)^{2}\sum_{i=1}^{n}a_{i}^{2},
\]

so \(\Delta=0\) and equality occurs.  
The trivial case \(\mathbf a=\mathbf 0\) also yields equality, because both sides of the inequality are zero.

Thus equality holds \emph{iff} the vectors \(\mathbf a\) and \(\mathbf b\) are proportional.

\end{proof}